\documentclass[12pt,a4paper]{article}
\usepackage[utf8]{inputenc}
\usepackage[spanish]{babel}
\usepackage{amsmath}
\usepackage{amsfonts}
\usepackage{amssymb}
\usepackage{graphicx}
\usepackage{geometry}
\usepackage{color}
\usepackage{hyperref}

\geometry{margin=2.5cm}

\title{\textbf{Reporte de Avance: Optimización de Rastreo y Análisis de Comportamiento en Ratones (EPM)}}
\author{Ingeniería en Sistemas - Tesis TT\_Ratones\_2026}
\date{5 de marzo de 2026}

\begin{document}

\maketitle

\section{Resumen del Avance}
En esta sesión de trabajo se logró consolidar una arquitectura de rastreo visual de estado del arte para el proyecto de tesis. La principal problemática era el "brinco" o inestabilidad del rastreador debido al \textit{domain shift} (pelaje albino sobre fondos claros) y a la oclusión en los brazos cerrados del laberinto en cruz elevado (EPM).

\section{Metodología y Resultados}

\subsection{Entrenamiento de YOLOv11 V3}
Se procedió a entrenar localmente un modelo \textbf{YOLOv11n} utilizando un dataset optimizado de \textbf{2,265 imágenes} anotadas específicamente para las condiciones del laboratorio (piso de granito y laberinto azul).
\begin{itemize}
    \item \textbf{Precisión (P):} 99.5\%
    \item \textbf{Recall (R):} 98.7\%
    \item \textbf{mAP50:} 99.4\%
\end{itemize}
El resultado es un modelo extremadamente robusto que ya no pierde al ratón ni confunde sombras con sujetos de estudio.

\subsection{Arquitectura del Tracker}
Se refinó el script \texttt{generar\_video\_prediccion.py} para implementar una división clara de tareas:
\begin{enumerate}
    \item \textbf{DLC (DeepLabCut):} Pose Estimation de extremidades y nariz.
    \item \textbf{SimBA:} Clasificación de comportamientos patológicos (Tigmotaxis).
    \item \textbf{YOLOv11:} Único responsable del rastreo visual (Punto Rojo) y conteo de tiempos por brazo.
\end{enumerate}

\section{Nuevas Herramientas Funcionales}

\subsection{Selector de ROIs del Investigador}
Se implementó un módulo interactivo que permite al investigador dibujar los rectángulos exactos del laberinto al inicio de cada proceso. Esto permite:
\begin{itemize}
    \item Adaptabilidad a cambios de perspectiva de la cámara.
    \item Conteo milimétrico de segundos en Brazo Norte, Sur, Este, Oeste y Centro.
    \item Opción de corrección (\textit{Undo}) durante la selección.
\end{itemize}

\subsection{Visualización HUD Avanzada}
El sistema de visualización (Head-Up Display) fue rediseñado bajo los siguientes criterios de usabilidad:
\begin{itemize}
    \item \textbf{Split HUD:} Información de comportamiento arriba a la derecha y tiempos de laberinto abajo a la derecha.
    \item \textbf{Opacidad al 90\%:} Maximiza la legibilidad de los datos científicos sobre el video.
    \item \textbf{Punto Dinámico:} El color del rastreador cambia automáticamente (Coral para brazos abiertos, Cyan para cerrados, Naranja para centro) para validación rápida.
\end{itemize}

\section{Resultados Visuales}
A continuación se presenta una captura del sistema en funcionamiento, donde se aprecia la nueva interfaz dividida y la detección precisa del sujeto experimental en el laberinto.

\begin{figure}[h]
    \centering
    \centering
    % RUTA DEL SCREENSHOT: reportes/screenshot_hud_final.png
    %\includegraphics[width=0.8\textwidth]{reportes/screenshot_hud_final.png}
    \caption{Visualización final del sistema con HUD dividido y tracking dinámico.}
\end{figure}

\section{Conclusión}
El sistema actual es completamente funcional, estable y listo para procesamiento por lotes masivo. Se ha eliminado el ruido en los datos de trayectorias, asegurando que las estadísticas finales de la tesis sean precisas y replicables.

\end{document}
